\chapter{Facelift}

\section*{Introduction \& Clinical Indications}
A facelift, formally known as rhytidectomy, is a complex cosmetic surgical procedure meticulously designed to address the visible signs of aging in the mid-face, lower face, and neck. It primarily targets sagging skin, deep facial folds such as nasolabial and marionette lines, the formation of jowls along the jawline, and loose skin with prominent platysmal bands in the neck. These changes are primarily driven by gravity, chronic sun exposure, genetic predisposition, and the natural process of collagen and elastin degradation. The procedure involves tightening the underlying facial muscles and the Superficial Musculoaponeurotic System (SMAS), repositioning descended fat pads, and redraping the skin to create a smoother, firmer, and more youthful contour. This intervention aims to restore a more refreshed and rejuvenated appearance, significantly improving the definition of the jawline and the cervicomental angle.

\begin{itemize}
  \item Rhytidectomy
  \item Facial rejuvenation surgery
  \item Cervicofacial rhytidectomy
  \item SMAS lift
  \item Deep plane facelift
  \item Composite facelift
\end{itemize}

The fundamental principle of a facelift is to counteract the multi-layered effects of gravitational descent, volume redistribution, and loss of skin and soft tissue elasticity that characterize facial aging. Modern facelift techniques move beyond simple skin tightening to address the deeper anatomical structures responsible for facial sagging, thereby achieving more natural and long-lasting results. The primary operative concept involves the precise manipulation and repositioning of the Superficial Musculoaponeurotic System (SMAS) and the platysma muscle.

By lifting, tightening, and often partially excising or plicating the SMAS layer, the surgeon can effectively restore the youthful position of descended cheek fat pads and muscle groups, re-establish a defined jawline, and smooth the cervicomental angle. This deeper tissue manipulation reduces tension on the overlying skin, minimizing the "pulled" or "wind-swept" appearance associated with older, skin-only techniques. The surgical rationale is to create a stable, elevated foundation for the facial soft tissues, allowing the skin to be redraped without excessive tension, leading to a more natural, harmonious, and durable rejuvenation. Technical goals include:
\begin{itemize}
  \item Elevating and securing descended malar (cheek) fat pads.
  \item Eliminating or significantly reducing jowls along the mandibular border.
  \item Tightening lax platysma muscles in the neck to smooth bands and improve the cervicomental angle.
  \item Removing excess skin without creating undue tension, ensuring inconspicuous scar placement.
  \item Restoring a youthful facial contour and volume distribution.
\end{itemize}

The concept of surgical facial rejuvenation dates back to the early 20th century, with the first documented facelift performed in 1901 by Eugen H\"oller in Germany, involving simple excision of skin in front of the ears. Early procedures, often termed "skin lifts," focused solely on removing excess skin, frequently leading to unnatural, stretched appearances and short-lived results due to the underlying tissues continuing to sag. Significant advancements occurred in the mid-20th century with surgeons like Jacques Joseph and Archibald McIndoe refining techniques, but the fundamental limitation of skin-only lifts persisted.

A pivotal moment in facelift history came in 1976 when Tord Skoog described the subplatysmal dissection, emphasizing the importance of the platysma muscle in neck rejuvenation. Concurrently, Mitz and Peyronie published their seminal work on the Superficial Musculoaponeurotic System (SMAS), identifying it as a distinct anatomical layer crucial for facial support. This discovery revolutionized facelift surgery by shifting the focus from merely tightening skin to addressing the deeper fascial layers. Subsequent innovations, including the high-SMAS facelift, the deep plane facelift by Sam Hamra in the 1990s, and the composite facelift, further refined the ability to lift and reposition facial tissues as a single, integrated unit. These advancements have led to more natural, long-lasting outcomes, reduced tension on the skin, and improved safety profiles, marking a continuous evolution towards comprehensive, multi-layered facial rejuvenation.

A facelift is indicated for individuals seeking significant and lasting improvement in the visible signs of facial and neck aging that are not adequately addressed by non-surgical modalities. Specific clinical indications include:
\begin{itemize}
  \item Significant skin laxity and sagging in the mid-face and lower face, leading to a tired or aged appearance.
  \item Deepening of the nasolabial folds, which are the lines extending from the nose to the corners of the mouth.
  \item Prominent marionette lines, which are the lines extending downwards from the corners of the mouth, giving a sad or stern expression.
  \item The presence of jowls, characterized by loose skin and fat that obscure the natural definition of the jawline.
  \item Loss of definition along the jawline and chin, often accompanied by a poorly defined cervicomental angle.
  \item Loose skin and excess fat under the chin and along the neck, frequently presenting as a "turkey wattle" deformity or visible platysmal bands.
  \item A desire for a more youthful, refreshed, and harmonious facial appearance that aligns with the patient's self-perception.
\end{itemize}

A facelift is predominantly a primary elective cosmetic procedure, meaning it is typically the first surgical intervention undertaken by a patient to address moderate to severe signs of facial and neck aging. It is not usually performed for medical necessity but rather to achieve aesthetic goals, offering the most comprehensive and durable solution for significant skin laxity and soft tissue descent. Non-surgical treatments such as dermal fillers, neurotoxins, and energy-based devices are often considered first-line for milder signs of aging or for patients who are not surgical candidates or prefer less invasive options.

However, a facelift can also serve as a secondary or adjunct procedure in several contexts. For instance, a revision facelift may be performed years after an initial facelift to address continued aging, to refine results, or to correct suboptimal outcomes from a previous surgery. Furthermore, a facelift is frequently combined with other facial rejuvenation procedures, such as blepharoplasty (eyelid surgery), brow lift, or laser skin resurfacing, to achieve a more comprehensive and harmonious overall facial aesthetic. In these cases, it acts as an adjunct, providing the foundational lift while other procedures address specific regional concerns. The decision to undergo a facelift, whether primary or secondary, is always based on a thorough consultation between the patient and surgeon, considering the patient's specific concerns, anatomical findings, and realistic expectations.

\section*{Relevant Surgical Anatomy}
A comprehensive understanding of facial and neck anatomy is paramount for safe and effective rhytidectomy. The face and neck are composed of several layers: the skin, subcutaneous fat, the Superficial Musculoaponeurotic System (SMAS), and deeper structures including muscles, nerves, and vessels. The SMAS is a crucial fibromuscular layer that covers the muscles of facial expression, is continuous with the platysma muscle in the neck inferiorly, and the temporoparietal fascia and galea aponeurotica superiorly. Manipulation of the SMAS is central to modern facelift techniques, providing durable lift and avoiding excessive skin tension.

Key anatomical structures and their surgical relevance include:
\begin{itemize}
  \item \textbf{Facial Nerve (CN VII):} Exiting the stylomastoid foramen, its five main branches (temporal, zygomatic, buccal, marginal mandibular, cervical) innervate the muscles of facial expression. These branches lie deep to the SMAS in most areas, making careful SMAS dissection and manipulation critical to avoid injury. The marginal mandibular nerve is particularly vulnerable as it crosses the facial artery and vein near the angle of the mandible.
  \item \textbf{Great Auricular Nerve:} A sensory nerve from the cervical plexus (C2-C3), it crosses the sternocleidomastoid muscle and supplies sensation to the earlobe and skin over the mastoid process. It is susceptible to injury during posterior skin flap dissection, leading to temporary or permanent numbness of the earlobe.
  \item \textbf{Platysma Muscle:} A broad, thin sheet of muscle in the neck, continuous with the SMAS. Laxity and separation of its medial borders contribute to prominent neck bands and the "turkey wattle" deformity, often requiring plication (platysmaplasty) during facelift.
  \item \textbf{Arterial Supply:} The skin flaps are primarily supplied by perforators from the facial artery, superficial temporal artery, and posterior auricular artery. Meticulous dissection and hemostasis are essential to preserve flap viability and prevent hematoma formation.
  \item \textbf{Surgical Landmarks:} Incisions are typically placed along the preauricular crease, within the tragus, around the earlobe, and into the postauricular sulcus and posterior hairline to camouflage scars. The mastoid process and angle of the mandible serve as important bony reference points during dissection and suspension.
\end{itemize}

\section*{Preoperative Workup \& Preparation}
Thorough preoperative workup and preparation are crucial for ensuring patient safety and optimizing the outcome of a facelift. This process begins with a comprehensive consultation and includes a detailed medical evaluation to identify any potential risks.

\begin{itemize}
  \item \textbf{Comprehensive Medical Evaluation:} A complete medical history and physical examination are performed to assess overall health, identify any pre-existing conditions (e.g., hypertension, diabetes, cardiac disease), and ensure the patient is a suitable candidate for surgery. This includes baseline blood tests (e.g., complete blood count, coagulation profile, electrolytes) and sometimes an electrocardiogram (ECG) or chest X-ray, particularly for older patients or those with comorbidities.
  \item \textbf{Anesthesia Clearance:} Patients with significant medical comorbidities may require clearance from their primary care physician or a specialist (e.g., cardiologist, pulmonologist) to ensure they are fit for general anesthesia.
  \item \textbf{Medication Review and Adjustment:} Patients must disclose all prescription medications, over-the-counter drugs, supplements, and herbal remedies. Blood-thinning medications (e.g., aspirin, NSAIDs, warfarin, clopidogrel, certain herbal supplements like ginkgo biloba, ginseng, or high-dose vitamin E) must be stopped typically 1-2 weeks prior to surgery to minimize bleeding risk, under the guidance of the prescribing physician.
  \item \textbf{Smoking Cessation:} Smoking significantly impairs wound healing, increases the risk of skin necrosis, and can lead to other complications. Patients are strongly advised to stop smoking and nicotine products at least 4-6 weeks before and after surgery.
  \item \textbf{NPO Status:} Patients are instructed to refrain from eating or drinking (NPO) for a specified period, usually 6-8 hours, before surgery to prevent aspiration during anesthesia.
  \item \textbf{Prophylactic Antibiotics:} A single dose of intravenous prophylactic antibiotics is typically administered preoperatively to reduce the risk of surgical site infection.
  \item \textbf{Informed Consent:} A detailed discussion of the procedure, its potential benefits, risks, alternatives, and expected recovery is conducted, and the patient provides informed consent.
  \item \textbf{Pre-operative Photography:} Standardized photographs are taken from multiple angles to document the patient's pre-operative appearance for comparison with post-operative results and for medical record purposes.
  \item \textbf{Logistical Planning:} Patients are advised to arrange for transportation home after surgery and for assistance during the initial recovery period.
\end{itemize}

A facelift is an elective cosmetic procedure, and careful patient selection is paramount to ensure safety and satisfactory outcomes.

\textbf{Situations requiring postponement, optimization, or an alternative plan:}
\begin{itemize}
  \item Severe, uncontrolled systemic medical conditions that pose an unacceptable anesthetic or surgical risk (e.g., unstable angina, recent myocardial infarction, uncontrolled hypertension, severe diabetes with end-organ damage, significant untreated coagulopathy).
  \item Active infection in the surgical area or systemic infection, which must be resolved prior to surgery.
  \item Severe body dysmorphic disorder or unrealistic expectations regarding the surgical outcome, which may lead to patient dissatisfaction despite a technically successful procedure.
  \item Inability or unwillingness to cease smoking for the required pre- and post-operative period, due to the significantly elevated risk of skin necrosis and wound healing complications.
  \item Pregnancy or breastfeeding.
\end{itemize}

\textbf{Relative Contraindications and Precautions:}
\begin{itemize}
  \item Heavy smoking or nicotine use, even if temporarily ceased, still carries a higher risk of complications.
  \item Significant bleeding disorders or chronic use of anticoagulants that cannot be safely discontinued for the perioperative period.
  \item Previous radiation therapy to the face or neck, which can compromise tissue vascularity, elasticity, and healing capacity.
  \item Certain autoimmune diseases or connective tissue disorders (e.g., lupus, scleroderma) that may impair wound healing or increase scar formation.
  \item Psychological instability, unmanaged mental health conditions, or a history of poor compliance with medical advice.
  \item Extreme skin laxity or poor skin elasticity, which may limit the extent of improvement and increase the risk of complications.
  \item History of keloid or hypertrophic scarring, requiring careful incision placement and postoperative scar management.
  \item Morbid obesity, which can increase anesthetic risks and technical challenges.
\end{itemize}

\section*{Surgical Technique \& Procedure}
A facelift is typically performed under general anesthesia or deep intravenous sedation, ensuring patient comfort and immobility throughout the procedure. The patient is positioned supine, often with the head slightly elevated to minimize venous congestion. The specific technique employed varies based on the patient's anatomy, the extent of aging, and desired outcome (e.g., SMAS plication, SMASectomy, deep plane lift), but common steps include:

\begin{enumerate}
  \item \textbf{Anesthesia and Positioning:} General anesthesia or deep sedation is administered. The patient is positioned supine with the head elevated, and the face and neck are prepped and draped in a sterile fashion.
  \item \textbf{Incision Placement:} Meticulously planned incisions are made to be inconspicuous. They typically begin in the temporal hairline, extend downwards in front of the ear (preauricular crease, often incorporating the tragus), loop around the earlobe, and continue behind the ear into the posterior hairline. A small submental incision (under the chin) may also be made for neck contouring.
  \item \textbf{Skin and Subcutaneous Dissection:} Through these incisions, the skin and subcutaneous fat are carefully elevated from the underlying Superficial Musculoaponeurotic System (SMAS) layer. The extent of dissection varies by technique, but it typically extends anteriorly to the nasolabial fold and inferiorly over the neck.
  \item \textbf{SMAS Manipulation:} The SMAS layer is then addressed. This may involve plication (folding and suturing), imbrication (overlapping and suturing), excision of a strip of SMAS (SMASectomy), or a deeper dissection plane (deep plane facelift) where the SMAS and overlying skin/fat are lifted as a composite unit. The goal is to tighten and reposition the SMAS superiorly and posteriorly to lift the mid-face and jawline.
  \item \textbf{Neck Contouring:} If indicated, liposuction of the neck and jowls is performed to remove excess fat. The platysma muscles in the neck are often tightened (platysmaplasty) through the submental incision to smooth neck bands and improve the cervicomental angle.
  \item \textbf{Hemostasis:} Meticulous hemostasis is achieved throughout the procedure using electrocautery to minimize bleeding and reduce the risk of postoperative hematoma.
  \item \textbf{Skin Redraping and Excision:} Once the underlying tissues are repositioned and secured, the excess skin is carefully trimmed, and the remaining skin is redraped smoothly over the new contours without tension.
  \item \textbf{Closure and Dressing:} The incisions are then closed with fine sutures or staples. Small drains may be placed temporarily under the skin to collect any accumulating fluid or blood. A compression dressing or garment is applied to minimize swelling and support the healing tissues. The entire procedure can take anywhere from 3 to 6 hours, depending on the complexity and extent of the lift.
\end{enumerate}

\section*{Complications \& Risks}
As with any surgical procedure, a facelift carries potential risks and complications, which are thoroughly discussed with the patient prior to surgery. These can be broadly categorized as general surgical risks and procedure-specific risks.

\textbf{General Surgical Risks:}
\begin{itemize}
  \item \textbf{Anesthesia Complications:} These can include nausea, vomiting, allergic reactions to medications, respiratory problems, or, rarely, more severe cardiac or neurological events such as stroke or myocardial infarction.
  \item \textbf{Bleeding/Hematoma:} Post-operative bleeding can lead to a hematoma (a collection of blood under the skin), which is the most common complication and may require surgical drainage. This can cause pain, swelling, and potentially compromise skin flap viability if not promptly addressed.
  \item \textbf{Infection:} Although rare due to prophylactic antibiotics, infection can occur at the incision sites or within the surgical field, potentially requiring antibiotic treatment or further surgical intervention.
  \item \textbf{Seroma:} A collection of clear fluid (seroma) can accumulate under the skin, sometimes requiring aspiration or drainage.
  \item \textbf{Deep Vein Thrombosis (DVT) / Pulmonary Embolism (PE):} Although rare in elective cosmetic surgery, blood clots can form in the legs and potentially travel to the lungs, a life-threatening complication. Prophylactic measures are often employed.
\end{itemize}

\textbf{Procedure-Specific Risks:}
\begin{itemize}
  \item \textbf{Facial Nerve Injury:} Temporary or, rarely, permanent injury to branches of the facial nerve (CN VII) can result in weakness or paralysis of facial muscles, leading to asymmetry of the face, difficulty smiling, or drooping of the lip or eyebrow. The marginal mandibular nerve is particularly vulnerable.
  \item \textbf{Skin Necrosis:} Impaired blood supply to the skin flaps can lead to skin death, particularly in smokers, resulting in open wounds, delayed healing, and potentially noticeable scarring.
  \item \textbf{Numbness or Altered Sensation:} Temporary or permanent numbness, tingling, or altered sensation in parts of the face or neck is common due to nerve manipulation, especially involving the great auricular nerve.
  \item \textbf{Poor Scarring:} While incisions are typically well-hidden, some individuals may develop hypertrophic scars (raised, red) or keloids (excessively overgrown scars), particularly in predisposed individuals.
  \item \textbf{Hair Loss (Alopecia):} Temporary or permanent hair loss can occur along the incision lines in the temporal or posterior hairline, which may require hair transplantation or scar revision.
  \item \textbf{Asymmetry:} Despite careful surgical technique, some degree of facial asymmetry may persist or develop postoperatively.
  \item \textbf{Unsatisfactory Cosmetic Outcome:} The patient may not be fully satisfied with the aesthetic result, which could be due to residual laxity, unnatural appearance, or other factors, potentially requiring revision surgery.
  \item \textbf{Pigmentary Changes:} Skin discoloration or hyperpigmentation can occur, especially in patients with darker skin tones.
\end{itemize}

\section*{Postoperative Care \& Recovery}
Immediately after the facelift, the patient is transferred to the Post-Anesthesia Care Unit (PACU) for close monitoring as they recover from anesthesia. Pain medication is administered as needed, and the head is kept elevated to help reduce swelling. A compression dressing or garment is typically applied to the face and neck, and small drains may be in place to prevent fluid accumulation. Patients are usually discharged home the same day or the following morning, depending on the extent of surgery and their recovery.

Over the first 24-48 hours, patients will experience swelling, bruising, and mild to moderate discomfort, which can be managed with prescribed oral pain medication. Drains, if used, are usually removed within 1-3 days. Patients are advised to keep their head elevated, avoid bending, lifting, or straining, and apply cold compresses to minimize swelling. Most visible sutures or staples are removed within 7-10 days. Significant bruising and swelling typically subside over the first 2-3 weeks, though residual swelling can persist for several months. Patients can usually return to light, non-strenuous activities and desk work after 2 weeks and resume more vigorous exercise, heavy lifting, and contact sports after 4-6 weeks. The final results of the facelift become progressively more apparent as swelling resolves and tissues settle, with scar maturation continuing for up to a year or more.

During a facelift, the surgeon documents several key findings that guide the operative strategy and contribute to the overall outcome. These typically include the degree of skin laxity, the extent of jowl formation, the prominence of platysmal bands in the neck, the descent of malar fat pads, and the integrity and thickness of the Superficial Musculoaponeurotic System (SMAS). The presence of excess subcutaneous fat in the neck and jowls is also noted, as this often requires liposuction. The surgeon's findings directly inform the choice of SMAS manipulation technique (e.g., plication, imbrication, deep plane lift) and the extent of skin excision.

In the context of a facelift, a traditional pathology report on resected tissues is generally not generated unless specific skin lesions (e.g., suspicious moles, basal cell carcinoma, squamous cell carcinoma) are incidentally excised with the skin flap. The primary resected tissue is excess skin and subcutaneous fat, which is typically discarded without formal pathological examination unless there is a clinical indication for microscopic analysis. Therefore, the "pathology" of a facelift is predominantly macroscopic, focusing on the anatomical changes observed and corrected during surgery, rather than microscopic tissue diagnosis. The success of the procedure is assessed by the aesthetic improvement and the naturalness of the rejuvenated appearance.

A normal and successful postoperative course following a facelift is characterized by a predictable progression of healing and resolution of initial symptoms, leading to a harmonious and natural-looking rejuvenation. Immediately after surgery, patients will experience mild to moderate pain, managed effectively with oral analgesics, along with swelling and bruising that are most pronounced in the first few days. Pain should gradually subside over the first week. The initial tightness and numbness in the face and neck are expected and will gradually improve over several weeks to months.

Within the first 2-3 weeks, the majority of visible bruising and significant swelling will resolve, allowing the patient to return to most social activities. The incisions will begin to heal, and scars will initially appear pink but will progressively fade and soften over 6-12 months, becoming inconspicuous. The jawline will appear more defined, jowls will be reduced, and the neck will be smoother. Full sensation may take several months to return, and some areas of permanent numbness are possible but usually well-tolerated. The final aesthetic results, including the natural contour and improved cervicomental angle, will become fully apparent as all residual swelling dissipates and tissues settle, typically by 3-6 months post-surgery. Patients should expect a refreshed and more youthful appearance, not a completely different face.

Postoperative care for a facelift involves meticulous wound management and a series of follow-up appointments to monitor healing, ensure optimal results, and address any concerns.

\begin{itemize}
  \item \textbf{Initial Post-operative Visits:} Typically scheduled at 1 week, 2-4 weeks, 3 months, 6 months, and 1 year after surgery. These visits allow the surgeon to assess wound healing, remove sutures or staples, monitor for complications, and evaluate the aesthetic outcome.
  \item \textbf{Suture/Staple Removal:} Usually performed within 7-14 days post-operatively, depending on the location and type of closure.
  \item \textbf{Wound Care:} Patients are instructed on gentle cleansing of incision lines and application of antibiotic ointment as directed. Sun protection is crucial for healing scars.
  \item \textbf{Activity Restrictions:} Patients are advised to gradually resume light activities after 2 weeks and avoid strenuous exercise, heavy lifting, and contact sports for 4-6 weeks to prevent increased swelling, bleeding, or wound dehiscence. Head elevation, especially during sleep, is recommended for several weeks.
  \item \textbf{Rehabilitation/Physical Therapy:} Not typically required, but gentle massage of healing areas may be recommended by the surgeon to aid in scar maturation and reduce persistent swelling or firmness.
  \item \textbf{Warning Signs:} Patients are instructed to contact their surgeon immediately if they experience fever (above 101$^{\circ}$F or 38.3$^{\circ}$C), severe or increasing pain unresponsive to medication, excessive or rapidly increasing swelling or bruising, pus or foul odor from incisions, new onset of facial weakness or numbness, or any signs of infection.
\end{itemize}
While a facelift provides long-lasting results, the natural aging process continues, and some individuals may consider minor touch-up procedures or non-surgical treatments years later to maintain their rejuvenated appearance.

\section*{Outcomes \& Prognosis}
Facelift surgery remains one of the most frequently performed cosmetic surgical procedures globally, consistently ranking among the top five aesthetic operations. According to statistics from major plastic surgery societies, hundreds of thousands of facelifts are performed annually worldwide, reflecting its enduring effectiveness in addressing the visible signs of facial aging. The procedure is predominantly sought by women, typically between the ages of 40 and 70, although a growing number of men are also undergoing facelifts, accounting for approximately 10-15\% of all cases. The primary indication for surgery is age-related changes, including significant skin laxity, prominent jowls, and neck sagging. While the mortality rate associated with facelifts is exceedingly low (estimated at less than 0.01\%), morbidity is primarily related to complications such as hematoma (incidence ranging from 1\% to 8\%, often requiring re-operation), facial nerve injury (temporary in 1-2\%, permanent in less than 0.5\%), and skin necrosis (especially in smokers, with rates up to 3-10\%). These complication rates underscore the importance of careful patient selection and meticulous surgical technique.

The outcomes of facelift surgery are generally highly favorable, with patient satisfaction rates typically ranging from 85\% to 95\%. The procedure effectively provides a significant aesthetic improvement, resulting in a more youthful, refreshed, and natural-looking appearance. Functional outcomes are primarily psychological, leading to enhanced self-esteem, improved body image, and increased confidence, with no direct impact on survival or physiological function. While a facelift can effectively "turn back the clock" by several years (often estimated at 5-10 years), it does not stop the natural aging process. Therefore, some degree of recurrent skin laxity and descent can be expected over time, typically becoming noticeable 5 to 10 years post-surgery, though patients will generally continue to look younger than if they had not undergone the procedure.

Prognostic factors influencing the longevity and quality of results include the patient's age, skin quality and elasticity, genetic predisposition to aging, lifestyle choices (e.g., smoking cessation, sun protection, healthy diet), and the specific surgical technique employed. Experienced surgeons utilizing advanced SMAS-based techniques (e.g., deep plane, high-SMAS) generally achieve more durable and natural-looking results compared to older, skin-only lifts. While there are no specific "landmark clinical trials" in the same vein as drug efficacy trials for facelifts, the evolution of techniques, particularly the introduction of SMAS concepts by Mitz and Peyronie and deep plane approaches by Hamra, represent pivotal advancements that have significantly improved long-term outcomes and are widely recognized as benchmarks in the field.

\section*{References}
\begin{enumerate}
  \item American Society of Plastic Surgeons. Facelift. Available from: \url{https://www.plasticsurgery.org/cosmetic-procedures/facelift} [Accessed 2024 May 15].
  \item Neligan PC, Gurtner GC, editors. Neligan's Plastic Surgery. 4th ed. Philadelphia, PA: Elsevier; 2020. Chapter 18: Facelift.
  \item Mitz V, Peyronie M. The superficial musculo-aponeurotic system (SMAS) in the parotid and cheek area. Plast Reconstr Surg. 1976 Jul;58(1):80-8.
  \item Hamra ST. The deep plane rhytidectomy. Plast Reconstr Surg. 1990 Jan;86(1):53-61.
\end{enumerate}

\clearpage
